% Chapter 7: Simulating paths.

\Cref{ch:06-monte-carlo} priced claims that depend on the stock only at maturity: one
normal draw per sample gives $S_T$ exactly. An Asian option needs the average of the
stock over a set of fixing dates, a barrier option needs to know whether the path crossed a
level, and a model without a closed-form solution, such as a local or stochastic volatility
model, cannot be sampled at $T$ in one step. In all three cases we must simulate a
\emph{path}: the vector $(S_{t_1},\dots,S_{t_m})$ on a grid of dates.

Two situations arise. When the joint law of the vector is known, as for geometric Brownian
motion with time-dependent parameters or for any Gaussian process, it can be sampled
exactly, with no error other than the statistical one, on any grid. Even then there is a
choice: the same vector can be built from the same normals in different orders, and the
order decides which normals carry most of the variance. This is irrelevant for plain Monte
Carlo and decisive for stratification and quasi-Monte Carlo. When the law is not known,
the stochastic differential equation must be discretised, and the scheme introduces a
bias. We measure the bias (strong and weak orders), remove its leading term (Richardson
extrapolation), and decide how to split a computing budget between finer time steps and
more paths (Duffie and Glynn).

Throughout, $0=t_0<t_1<\dots<t_m=T$ is a time grid with steps $\Delta_j=t_j-t_{j-1}$,
$W$ is a Brownian motion under the risk-neutral measure $\Q$, and $Z_1,Z_2,\dots$ are
independent $\mathcal N(0,1)$ variables. Path arrays in the code have one row per path.

\section{Exact simulation on a time grid}\label{sec:ps-exact}

\subsection{Geometric Brownian motion with term structures}

In the standard model of \cref{def:sm-model} the parameters are constants. Market data
are not: the rate and the yield are read off a curve, and implied volatilities depend on
the maturity. The simplest model consistent with such data lets $r$, $y$ and $\sigma$
depend on time. Let $r,y,\sigma:[0,T]\to\R$ be bounded Borel functions, $\sigma\ge0$, and
\begin{equation}\label{eq:ps-gbm-ts}
  \frac{\dd S_t}{S_t}=\big(r(t)-y(t)\big)\,\dd t+\sigma(t)\,\dd W_t,\qquad S_0>0 .
\end{equation}
For $0\le s<t$ write
\[
  \bar r(s,t)=\frac1{t-s}\int_s^tr(u)\,\dd u,\qquad
  \bar\sigma^2(s,t)=\frac1{t-s}\int_s^t\sigma(u)^2\,\dd u,
\]
and $\bar y(s,t)$ likewise. The \emph{term rate} and the \emph{term volatility} are
$r_{0,t}=\bar r(0,t)$ and $\sigma_{0,t}=\bar\sigma(0,t)$; the \emph{interval forward}
quantities on the grid are $\bar r_j=\bar r(t_{j-1},t_j)$, $\bar y_j$ and
$\bar\sigma_j^2=\bar\sigma^2(t_{j-1},t_j)$, and we set
$\bar\nu_j=\bar r_j-\bar y_j-\tfrac12\bar\sigma_j^2$. Forward variances are recovered from
term volatilities by
\begin{equation}\label{eq:ps-forward-variance}
  \bar\sigma_j^2\,\Delta_j=t_j\,\sigma_{0,t_j}^2-t_{j-1}\,\sigma_{0,t_{j-1}}^2 ,
\end{equation}
so a term structure of volatility is consistent with \eqref{eq:ps-gbm-ts} only if the
total variance $t\,\sigma_{0,t}^2$ is nondecreasing in $t$ (\exref{ex:ps-forward-variance}).

The integrals of a deterministic integrand against $W$ are Gaussian. We record the fact
with a proof, since it is the whole content of exact simulation.

\begin{lemma}[Wiener integrals]\label{lem:ps-wiener}
Let $g:[0,T]\to\R$ be a deterministic function with $\int_0^Tg^2<\infty$, and
$I_j=\int_{t_{j-1}}^{t_j}g(u)\,\dd W_u$ for $j=1,\dots,m$. Then $I_1,\dots,I_m$ are
independent and $I_j\sim\mathcal N\big(0,\int_{t_{j-1}}^{t_j}g^2\big)$.
\end{lemma}

\begin{proof}
If $g$ is a step function, constant on the intervals of a partition that refines the grid,
each $I_j$ is a finite
linear combination of increments of $W$ over disjoint subintervals of $(t_{j-1},t_j]$. So
$(I_1,\dots,I_m)$ is a linear image of independent Gaussian increments; distinct $I_j$ use
disjoint sets of increments, hence are independent, and $\Var(I_j)$ is the sum of
$g^2$ times the lengths of the subintervals, that is $\int_{t_{j-1}}^{t_j}g^2$.

For general $g$, pick step functions $g_n\to g$ in $L^2[0,T]$ (they are dense) and let
$I^{(n)}_j$ be the corresponding integrals. By It\^o's isometry,
$\E[(I^{(n)}_j-I_j)^2]=\int_{t_{j-1}}^{t_j}(g_n-g)^2\to0$. For $u\in\R^m$, using
$|e^{ia}-e^{ib}|\le|a-b|$ and the Cauchy--Schwarz inequality,
\[
  \Big|\E\big[e^{i\sum_ju_jI^{(n)}_j}\big]-\E\big[e^{i\sum_ju_jI_j}\big]\Big|
  \le\sum_j|u_j|\,\E\big[(I^{(n)}_j-I_j)^2\big]^{1/2}\longrightarrow0 .
\]
The left characteristic function equals $\prod_j\exp\big(-\tfrac12u_j^2\int_{t_{j-1}}^{t_j}g_n^2\big)$
by the step-function case, and converges to $\prod_j\exp\big(-\tfrac12u_j^2\int_{t_{j-1}}^{t_j}g^2\big)$.
A characteristic function determines the law, also in $\R^m$
\citep[Section~3.10]{durrett}, so the components are independent Gaussians.
\end{proof}

\begin{proposition}[Exact simulation of GBM]\label{prop:ps-gbm-exact}
\begin{enumerate}[label=(\roman*)]
  \item The unique solution of \eqref{eq:ps-gbm-ts} is
    $S_t=S_0\exp\big(\int_0^t\nu(u)\,\dd u+\int_0^t\sigma(u)\,\dd W_u\big)$, with
    $\nu=r-y-\tfrac12\sigma^2$.
  \item For every grid, the recursion
    \begin{equation}\label{eq:ps-gbm-step}
      \hat S_0=S_0,\qquad
      \hat S_j=\hat S_{j-1}\exp\big(\bar\nu_j\Delta_j+\bar\sigma_j\sqrt{\Delta_j}\,Z_j\big),
      \quad j=1,\dots,m,
    \end{equation}
    produces a vector $(\hat S_0,\dots,\hat S_m)$ with the same joint law as
    $(S_{t_0},\dots,S_{t_m})$, whatever the step sizes.
  \item In particular $\EQ[S_t]=S_0e^{(r_{0,t}-y_{0,t})t}$, and the European call with
    maturity $t$ has the Black--Scholes price with the term parameters $r_{0,t}$,
    $y_{0,t}$, $\sigma_{0,t}$ (if $\sigma_{0,t}=0$, the formula is read as its limit
    $\sigma\downarrow0$, the discounted intrinsic value of the forward).
\end{enumerate}
\end{proposition}

\begin{proof}
(i) Let $X_t=\int_0^t\nu+\int_0^t\sigma\,\dd W$, an It\^o process with
$\dd X=\nu\,\dd t+\sigma\,\dd W$ and $\dd[X,X]=\sigma^2\dd t$. It\^o's formula applied to
$S_0e^{X_t}$ gives $\dd S=S\,(\nu+\tfrac12\sigma^2)\,\dd t+S\sigma\,\dd W$, which is
\eqref{eq:ps-gbm-ts}. Its coefficients are Lipschitz in $S$ uniformly in $t$, because
$r,y,\sigma$ are bounded, and for such coefficients strong (pathwise) uniqueness holds
\citep[Theorem~5.2.5]{karatzas-shreve}.

(ii) By (i), $\log S_{t_j}-\log S_{t_{j-1}}=\bar\nu_j\Delta_j+I_j$ with
$I_j=\int_{t_{j-1}}^{t_j}\sigma\,\dd W$. By \cref{lem:ps-wiener} the $I_j$ are independent
with $I_j\sim\mathcal N(0,\bar\sigma_j^2\Delta_j)$, so the vector of log-increments has the
law of $\big(\bar\nu_j\Delta_j+\bar\sigma_j\sqrt{\Delta_j}Z_j\big)_{j\le m}$. Both
$(S_{t_j})_j$ and $(\hat S_j)_j$ are the same function (cumulative sums, then the
exponential) of these vectors.

(iii) $\log S_t$ is normal with mean $\log S_0+\int_0^t\nu$ and variance
$\int_0^t\sigma^2=\sigma_{0,t}^2t$, so
$\EQ[S_t]=S_0\exp\big(\int_0^t(r-y)\big)$. A lognormal $S_t$ with this mean and this
log-variance is exactly the terminal law in the Black--Scholes model with constant
parameters $r_{0,t}$, $y_{0,t}$, $\sigma_{0,t}$, and the discount factor
$e^{-\int_0^tr}=e^{-r_{0,t}t}$ is the same.
\end{proof}

Three remarks. The recursion has no discretisation error: a grid is chosen only because
the payoff looks at those dates. The parameters enter only through their averages over the
steps; this is the precise meaning of the ``effective'' parameters of a time-stepping
scheme. And nothing is said about the path between grid dates, which \cref{sec:ps-brownian}
fills in.

\begin{example}[A term structure from the course]\label{exa:ps-term}
Take $y=0$, an instantaneous forward rate $r(t)=0.02\,t$, a volatility $\sigma(t)=20\%$
on $[0,\frac12)$ and $10\%$ afterwards, and monthly dates $t_j=j/12$. The interval forward
variances are $\bar\sigma_j^2=0.04$ for $j\le6$ and $0.01$ for $j>6$; the interval forward
rates are $\bar r_j=0.01\,(t_{j-1}+t_j)=0.01\,(2j-1)/12$. At one year the term rate is
$r_{0,1}=1\%$ and the term volatility $\sigma_{0,1}=\sqrt{0.025}=15.81\%$, so the one-year
at-the-money call is worth $6.7799$ with $S_0=K=100$ (\cref{prop:ps-gbm-exact}(iii)), and
\eqref{eq:ps-gbm-step} reproduces this within its standard error
(\coderef{code/tests/test_paths.py}).
\end{example}

\begin{remark}[Frozen parameters]\label{rem:ps-frozen}
A tempting shortcut is to step with the parameters at the left end of each step,
$r(t_{j-1})$ and $\sigma(t_{j-1})$: this is the Euler scheme of \cref{sec:ps-schemes}
applied to $\log S$. In \cref{exa:ps-term} it replaces $\int_0^1r=0.01$ by the Riemann sum
$0.00917$, so the simulated forward is $100.92$ instead of $101.01$. Worse, if the
convention $\sigma(\frac12)=20\%$ is used at the jump, the variance $\int_0^1\sigma^2$
becomes $0.0275$ instead of $0.025$, and the call, discounted with the true factor
$e^{-0.01}$, comes out at $7.0384$ instead of $6.7799$, a $4\%$ error that no number of paths
will remove (with the convention $\sigma(\frac12)=10\%$ the drift error alone gives
$6.7337$; discounting with the Riemann sum as well gives $7.0443$ and $6.7393$). With averages
over the steps, the convention at the jump does not matter.
\end{remark}

\begin{remark}[Time change]\label{rem:ps-time-change}
Let $\Theta(t)=\int_0^t\sigma^2$. The process $M_t=\int_0^t\sigma\,\dd W$ is centered
Gaussian with $\Cov(M_s,M_t)=\Theta(s\wedge t)$ (\cref{lem:ps-wiener}), which is the
covariance of $B_{\Theta(t)}$ for a Brownian motion $B$. So on any grid, $M$ has the law of
a standard Brownian motion run on the \emph{variance clock} $\Theta$; the code builds GBM
paths this way. If moreover the drift is proportional to the variance rate,
$r(t)-y(t)=\lambda\sigma(t)^2$, then $\log(S_t/S_0)$ has the law of
$(\lambda-\frac12)\Theta(t)+B_{\Theta(t)}$, a Brownian motion with constant drift in the
new clock, and constant-coefficient formulas, such as those for barriers in
\cref{ch:16-barriers}, apply in that clock. In \cref{exa:ps-term}, $\Theta(1)=0.025$ and
$\int_0^1r=0.01$ give $\lambda=0.4$: a process with the same total drift and
variance but drift proportional to variance, of the kind used as a control variate for
discretely monitored barriers, has monthly forward rates $0.4\times0.04=1.6\%$ for the first
six months and $0.4\times0.01=0.4\%$ afterwards.
\end{remark}

\subsection{Gaussian processes}

\Cref{prop:ps-gbm-exact} is a case of a general recipe. If $X$ is a Gaussian process, the
vector $(X_{t_1},\dots,X_{t_m})$ is $\mathcal N(\mu,C)$ with $\mu_i=\E X_{t_i}$ and
$C_{ij}=\Cov(X_{t_i},X_{t_j})$, and $\mu+AZ$ has this law for any matrix $A$ with
$AA^\top=C$ (\cref{prop:rng-cholesky}). A dense factorisation costs $O(m^3)$ once and
$O(m^2)$ per path. When $X$ is also Markov the cost drops to $O(m)$ per path, because
each $X_{t_j}$ given the past is Gaussian with a mean that depends only on $X_{t_{j-1}}$.

\begin{example}[Ornstein--Uhlenbeck]\label{exa:ps-ou}
Let $\dd X_t=\kappa(\theta-X_t)\,\dd t+\xi\,\dd W_t$ with $\kappa>0$, the Vasicek short-rate
model. Applying It\^o's formula to $e^{\kappa t}X_t$ gives
$\dd(e^{\kappa t}X_t)=e^{\kappa t}(\kappa\theta\,\dd t+\xi\,\dd W_t)$, hence
\[
  X_{t+h}=\theta+(X_t-\theta)e^{-\kappa h}+\xi\int_t^{t+h}e^{-\kappa(t+h-u)}\,\dd W_u .
\]
By \cref{lem:ps-wiener} (and since the integral is independent of $\F_t$), the exact step is
$X_{t+h}=\theta+(X_t-\theta)e^{-\kappa h}+\xi\sqrt{(1-e^{-2\kappa h})/(2\kappa)}\,Z$, on
any grid. The Euler scheme for the same equation is biased, and unstable for $\kappa h>2$
(\exref{ex:ps-ou-euler}).
\end{example}

\demo{07-path-simulation.ipynb}{move the volatility jump and the number of steps, and compare
exact stepping with frozen parameters against the closed-form price.}

\begin{incode}
\coderef{code/src/mfe230d/paths.py}: \texttt{gbm\_paths(S0, t, r, y, sigma, n\_paths, seed)}
implements \eqref{eq:ps-gbm-step}; \texttt{r}, \texttt{y}, \texttt{sigma} may be numbers,
\texttt{PiecewiseConstant} curves or functions of time, and \texttt{effective\_parameters}
returns the averages $\bar r_j$, $\bar y_j$, $\bar\sigma^2_j$ and the clock $\Theta(t_j)$.
\texttt{ou\_paths} is the exact Ornstein--Uhlenbeck step. Every sampler accepts the normals
\texttt{Z} instead of a seed, so that antithetic, stratified or quasi-random normals can be
passed in (\cref{ch:08-antithetic-control,ch:10-stratification,ch:11-quasi-monte-carlo}).
\end{incode}

\section{Three ways to build a Brownian path}\label{sec:ps-brownian}

By \cref{rem:ps-time-change}, the Gaussian part of a GBM path is a Brownian motion sampled
on a grid (of the variance clock). Its covariance is $C_{ij}=\min(t_i,t_j)$, and each square root
$A$ of $C$ gives a construction $W=AZ$. Three are in use.

\emph{Increments.} $W_{t_j}=W_{t_{j-1}}+\sqrt{\Delta_j}\,Z_j$. The matrix $A$ is lower
triangular with $A_{ij}=\sqrt{\Delta_j}$ for $j\le i$: the Cholesky factor of $C$. The
normal $Z_j$ moves the path from $t_j$ on.

\emph{Brownian bridge.} Sample the endpoint first, $W_T=\sqrt T\,Z_1$, then the midpoint
given the two ends, then the quarter points given their neighbours, and so on. This rests
on the law of a Brownian path between two known values.

\begin{proposition}[Brownian bridge]\label{prop:ps-bridge}
Let $0\le s<u<t$ and let $\mathcal G_{s,t}=\sigma(W_r:r\in[0,s]\cup[t,\infty))$.
\begin{enumerate}[label=(\roman*)]
  \item The variable
    \[
      B=W_u-\frac{(t-u)\,W_s+(u-s)\,W_t}{t-s}
    \]
    is independent of $\mathcal G_{s,t}$ and $B\sim\mathcal N\big(0,\frac{(u-s)(t-u)}{t-s}\big)$.
  \item Consequently, given $\mathcal G_{s,t}$, $W_u$ is normal with mean
    $\frac{(t-u)W_s+(u-s)W_t}{t-s}$ (linear interpolation) and variance
    $\frac{(u-s)(t-u)}{t-s}$: it depends on the rest of the path only through the two
    neighbouring values.
  \item Let $0<t_1<\dots<t_m$ and fill the points in any order. When $t_i$ is filled, let
    $t_l$ be the largest already-filled time below $t_i$ (or $0$, with $W_0=0$) and $t_r$
    the smallest already-filled time above it. Set
    $W_{t_i}=W_{t_l}+\sqrt{t_i-t_l}\,Z$ if there is no $t_r$, and otherwise
    \begin{equation}\label{eq:ps-bridge-step}
      W_{t_i}=\frac{(t_r-t_i)\,W_{t_l}+(t_i-t_l)\,W_{t_r}}{t_r-t_l}
      +\sqrt{\frac{(t_i-t_l)(t_r-t_i)}{t_r-t_l}}\;Z ,
    \end{equation}
    with a fresh normal $Z$ at each step. The resulting vector has the law of
    $(W_{t_1},\dots,W_{t_m})$.
\end{enumerate}
\end{proposition}

\begin{proof}
(i) Write $a=u-s$, $b=t-u$, $D_1=W_u-W_s$ and $D_2=W_t-W_u$. Then
$B=W_s+D_1-W_s-\frac{a}{a+b}(D_1+D_2)=\frac{bD_1-aD_2}{a+b}$. The increments $D_1$, $D_2$
are independent with variances $a$, $b$, so $B$ is centered normal with variance
$\frac{b^2a+a^2b}{(a+b)^2}=\frac{ab}{a+b}$, and
$\Cov(B,D_1+D_2)=\frac{ba-ab}{a+b}=0$. As $(B,D_1+D_2)$ is Gaussian, $B$ is independent
of $D_1+D_2=W_t-W_s$.

Let $\mathcal H$ be generated by $(W_r)_{r\le s}$ and $(W_r-W_t)_{r\ge t}$. The pair
$(D_1,D_2)$ is independent of $\mathcal H$ (independent increments), and $\mathcal G_{s,t}=\sigma(\mathcal H,W_t-W_s)$ because
$W_r=W_s+(W_t-W_s)+(W_r-W_t)$ for $r\ge t$. For bounded measurable $f,g,k$,
\[
  \E\big[f(B)\,g(H)\,k(W_t-W_s)\big]=\E[g(H)]\,\E[f(B)]\,\E[k(W_t-W_s)]
  =\E[f(B)]\,\E\big[g(H)\,k(W_t-W_s)\big],
\]
where $H$ is any finite vector of $\mathcal H$-measurable variables: the first equality
uses that $(B,W_t-W_s)$ is a function of $(D_1,D_2)$, independent of $H$, and then the
independence of $B$ and $W_t-W_s$; the second uses again that $W_t-W_s$ is independent of
$H$. Such products determine independence from $\mathcal G_{s,t}$, by Dynkin's
$\pi$--$\lambda$ theorem \citep[Section~2.1]{durrett}.

(ii) $W_u$ is the sum of the $\mathcal G_{s,t}$-measurable interpolation and the
independent variable $B$; condition with the independence lemma
\citep[Lemma~2.3.4]{shreve2}.

(iii) By induction on the number $k$ of filled points: the joint law of the first $k$
values is that of the corresponding $W$'s. The $k$-th point $t_i$ is generated from its
filled neighbours by \eqref{eq:ps-bridge-step}. All previously filled times lie outside
$(t_l,t_r)$, so the $\sigma$-algebra $\mathcal F'$ they generate is contained in
$\mathcal G_{t_l,t_r}$ and contains $W_{t_l}$, $W_{t_r}$. By (ii) and the tower property, the
conditional law of $W_{t_i}$ given $\mathcal F'$ is the normal law in
\eqref{eq:ps-bridge-step}; without a right neighbour it is $\mathcal N(W_{t_l},t_i-t_l)$ by
independent increments. A joint law is determined by the law of the first $k-1$
coordinates and the conditional law of the $k$-th given them.
\end{proof}

The conditional variance $\frac{(u-s)(t-u)}{t-s}$ vanishes at both ends and is largest,
$\frac{t-s}4$, at the midpoint. The \emph{pinned} bridge $W_u-\frac uTW_T$, $u\in[0,T]$,
is independent of $W_T$ with covariance $\min(u,v)-uv/T$. On a grid of $m=2^k$ steps the
dyadic order (\texttt{bridge\_order} in the code) fills $T$, $T/2$, $T/4$ and $3T/4$, and
so on; after $2^\ell$ normals the path is known exactly at the dyadic points of level
$\ell$ (\cref{fig:ps-bridge}, left). The construction costs $O(m)$ per path, like the
increments, and works for any grid, uniform or not.

\emph{Principal components.} Diagonalise $C=V\Lambda V^\top$ with eigenvalues
$\lambda_1\ge\dots\ge\lambda_m>0$ and orthonormal eigenvectors $v_i$, and take
$A=V\Lambda^{1/2}$: $W=\sum_i\sqrt{\lambda_i}\,v_iZ_i$. This costs $O(m^2)$ per path.

All three constructions produce exactly the same law, so for plain Monte Carlo they give
estimators with the same variance; only their cost differs. They differ in how the
variance is spread over the normals. Call
\[
  \pi_k(A)=\frac{\sum_{j\le k}\|a_j\|^2}{\operatorname{tr}C},\qquad a_j\text{ the $j$-th column of }A,
\]
the share of the total variance $\sum_i\Var(W_{t_i})=\operatorname{tr}C$ carried by the
first $k$ normals.

\begin{proposition}[Principal components carry the most variance]\label{prop:ps-pca}
For every $m\times m$ matrix $A$ with $AA^\top=C$ and every $k\le m$,
$\pi_k(A)\le\sum_{i\le k}\lambda_i/\operatorname{tr}C$, with equality for $A=V\Lambda^{1/2}$.
\end{proposition}

\begin{proof}
$C$ is invertible (the increments over the grid are independent and nondegenerate), and
$U=C^{-1/2}A$ satisfies $UU^\top=C^{-1/2}CC^{-1/2}=I$, so its columns $u_j$ are
orthonormal and $a_j=C^{1/2}u_j$. Then
$\sum_{j\le k}\|a_j\|^2=\sum_{j\le k}u_j^\top Cu_j=\sum_i\lambda_iw_i$ with
$w_i=\sum_{j\le k}(v_i^\top u_j)^2$. Each $w_i$ is the squared length of the projection of
$v_i$ on the span of $u_1,\dots,u_k$, so $0\le w_i\le1$, and $\sum_iw_i=k$. A weighted sum
$\sum_i\lambda_iw_i$ with these constraints is largest when the weight $1$ sits on the $k$
largest eigenvalues, which gives the bound. For $A=V\Lambda^{1/2}$,
$\|a_j\|^2=\lambda_j$.
\end{proof}

For a uniform grid with $m=64$ steps, the shares are:
\begin{center}
\begin{tabular}{lcccc}
\toprule
first $k$ normals & $k=1$ & $k=2$ & $k=4$ & $k=16$\\
\midrule
increments & $0.031$ & $0.061$ & $0.120$ & $0.435$\\
Brownian bridge & $0.672$ & $0.836$ & $0.918$ & $0.981$\\
principal components & $0.811$ & $0.901$ & $0.950$ & $0.988$\\
\bottomrule
\end{tabular}
\end{center}
The first increment carries $2/(m+1)$ of the variance; the first bridge normal, which sets
$W_T$, carries $\sum_it_i^2/(T\sum_it_i)=(2m+1)/(3m)\to\frac23$; the first principal
component tends to $8/\pi^2\approx0.81$ (\exref{ex:ps-pca-eigen}). This ordering matters
whenever the normals are not treated symmetrically: stratifying $Z_1$ stratifies $W_T$
under the bridge construction (\cref{prop:st-bridge}), and quasi-Monte Carlo points are
most uniform in their first coordinates, which is why the bridge and PCA lower the
effective dimension of path-dependent integrands (\cref{rem:qmc-effective-dimension}). The
bridge law (ii) also answers questions about the path between two grid dates; the
probability that it crosses a barrier is \cref{prop:bar-bridge-crossing}.

\begin{figure}[ht]
  \centering
  \includegraphics[width=\textwidth]{ch07-bridge}
  \caption{Left: one Brownian path built by the bridge construction on $2^8$ steps, shown
    after $1$, $2$, $4$, $16$ and $256$ normals (linear interpolation between the points
    already built). Right: share $\pi_k$ of the variance carried by the first $k$ normals,
    $m=64$ (\cref{prop:ps-pca}).}
  \label{fig:ps-bridge}
\end{figure}

\demo{07-path-simulation.ipynb}{watch the bridge construction refine a path level by level,
and compare the variance carried by the first normals under the three constructions.}

\section{Discretising a stochastic differential equation}\label{sec:ps-schemes}

Now let $X$ solve
\begin{equation}\label{eq:ps-sde}
  \dd X_t=a(t,X_t)\,\dd t+b(t,X_t)\,\dd W_t,\qquad X_0=x_0\in\R,
\end{equation}
for which no exact sampler is available. Take the uniform grid $t_k=kh$, $h=T/n$, and
write $\Delta W_k=W_{t_{k+1}}-W_{t_k}$. Over one step,
\[
  X_{t_{k+1}}=X_{t_k}+\int_{t_k}^{t_{k+1}}a(s,X_s)\,\dd s+\int_{t_k}^{t_{k+1}}b(s,X_s)\,\dd W_s .
\]
Freezing both integrands at the left end gives the \emph{Euler} (Euler--Maruyama) scheme
\begin{equation}\label{eq:ps-euler}
  \bar X_0=x_0,\qquad \bar X_{k+1}=\bar X_k+a(t_k,\bar X_k)\,h+b(t_k,\bar X_k)\,\Delta W_k .
\end{equation}
The largest term it neglects comes from the variation of $b(s,X_s)$ inside the step. By
It\^o's formula, $b(s,X_s)=b(t_k,X_{t_k})+\int_{t_k}^s(b\,\partial_xb)(r,X_r)\,\dd W_r+{}$a
$\dd r$-integral, and freezing $b\,\partial_xb$ at $t_k$ leaves
$\int_{t_k}^{t_{k+1}}(W_s-W_{t_k})\,\dd W_s=\frac12\big((\Delta W_k)^2-h\big)$ (It\^o's
formula for $(W_s-W_{t_k})^2$). This term has standard deviation $h/\sqrt2$, while the
other neglected terms are of order $h^{3/2}$. Keeping it gives the \emph{Milstein} scheme
\begin{equation}\label{eq:ps-milstein}
  \tilde X_{k+1}=\tilde X_k+a(t_k,\tilde X_k)\,h+b(t_k,\tilde X_k)\,\Delta W_k
  +\tfrac12\,(b\,\partial_xb)(t_k,\tilde X_k)\big((\Delta W_k)^2-h\big).
\end{equation}
When $b$ does not depend on $x$ (additive noise) the two schemes coincide.

\begin{definition}[Strong and weak order]\label{def:ps-orders}
A scheme $Y^h$ with step $h$ has \emph{strong order} $\gamma>0$ if
$\E|X_T-Y^h_T|\le C\,h^\gamma$ for all $h\le h_0$, and \emph{weak order} $\beta>0$ for a
class of functions $f$ if $|\E f(Y^h_T)-\E f(X_T)|\le C_f\,h^\beta$ for all $h\le h_0$.
\end{definition}

The strong error compares the two paths driven by the same Brownian motion; it is what
matters when the path itself is used, as in hedging simulations or the couplings of
multilevel Monte Carlo. The weak error compares laws, and is the bias of a Monte Carlo
price. Since $|\E f(Y^h_T)-\E f(X_T)|\le L\,\E|Y^h_T-X_T|$ for an $L$-Lipschitz $f$, strong
order $\gamma$ implies weak order $\gamma$ for such $f$; the weak order is often higher.

\begin{theorem}[Euler--Maruyama]\label{thm:ps-euler}
Suppose there is $K$ such that, for all $s,t\in[0,T]$ and $x,z\in\R$,
\begin{enumerate}[label=(A\arabic*)]
  \item $|a(t,x)-a(t,z)|+|b(t,x)-b(t,z)|\le K|x-z|$;
  \item $|a(t,x)|+|b(t,x)|\le K(1+|x|)$;
  \item $|a(s,x)-a(t,x)|+|b(s,x)-b(t,x)|\le K(1+|x|)\,|s-t|^{1/2}$.
\end{enumerate}
\begin{enumerate}[label=(\roman*)]
  \item The Euler scheme has strong order $\frac12$: $\E|X_T-\bar X^h_T|\le C\,h^{1/2}$
    \citep[Theorem~10.2.2]{kloeden-platen}.
  \item Suppose moreover that $a$ and $b$ are infinitely differentiable on $[0,T]\times\R$
    with bounded partial derivatives of every order $\ge1$, and that $f$ is infinitely
    differentiable with $f$ and all its derivatives of polynomial growth. Then
    \begin{equation}\label{eq:ps-tt}
      \E f(\bar X^h_T)-\E f(X_T)=C_f\,h+O(h^2)
    \end{equation}
    for a constant $C_f$ \citep{talay-tubaro}; in particular the weak order is $1$.
    Weak order $1$ holds already when $a$, $b$, $f$ have four derivatives of polynomial
    growth \citep[Theorem~14.5.1 with $\beta=1$]{kloeden-platen}.
\end{enumerate}
\end{theorem}

\Citet{talay-tubaro} state (ii) for time-independent coefficients in several dimensions;
the pair $(t,X_t)$ solves such an equation, and its Euler scheme is \eqref{eq:ps-euler}. The
proof of (i) compares the scheme with the solution in mean square and closes with
Gronwall's lemma. The proof of (ii) teaches more: it explains why the weak order of Euler is
twice its strong order.

\begin{proof}[Proof of (ii), sketch]
Let $u(t,x)=\E[f(X_T)\mid X_t=x]$. Under the hypotheses, $u$ is smooth with derivatives of
polynomial growth \citep{talay-tubaro}, and by the Markov property
$\E u(t+h,X^{t,x}_{t+h})=u(t,x)$ for the solution $X^{t,x}$ started at $x$ at time $t$.
Telescoping,
\begin{gather*}
  \E f(\bar X_n)-\E f(X_T)=\E u(t_n,\bar X_n)-u(0,x_0)=\sum_{k=0}^{n-1}\E\,\varepsilon(t_k,\bar X_k),\\
  \varepsilon(t,x)=\E\big[u(t+h,x+\hat D)-u(t+h,x+D)\big],
\end{gather*}
where $\hat D=a(t,x)h+b(t,x)\Delta W$ is one Euler step and $D=X^{t,x}_{t+h}-x$ one exact
step. Expand $u(t+h,\cdot)$ around $x$ to third order: $\varepsilon$ is a combination of the
differences $\E\hat D^j-\E D^j$, $j=1,2,3$, plus a remainder of order $(1+|x|^q)h^2$,
because $\partial_x^4u$ grows polynomially and both steps have fourth moments of order
$h^2$. For the Euler step,
$\E\hat D=ah$, $\E\hat D^2=b^2h+a^2h^2$, $\E\hat D^3=3ab^2h^2+a^3h^3$, $\E\hat D^4=O(h^2)$.
For the exact step, It\^o's formula applied to $(X_s-x)^j$ gives $\E D=ah+O(h^2)$,
$\E D^2=b^2h+O(h^2)$, $\E D^3=O(h^2)$, $\E D^4=O(h^2)$. So
$|\varepsilon(t,x)|\le C(1+|x|^q)h^2$. The scheme has moments of all orders bounded uniformly
in $h$ and $k$ (a discrete Gronwall argument using (A2)), and summing $n=T/h$ local errors of
order $h^2$ gives an error of order $h$.
\end{proof}

The local errors of Euler have size $h$ in mean square but $h^2$ in mean. Over $n=T/h$ steps,
errors with zero mean add up in quadrature, to $\sqrt n\,h=O(h^{1/2})$: the strong order. Their
means add up linearly, to $n\,h^2=O(h)$: the weak order. The proof also shows that only the
first three moments of the increments matter for the weak error, so $\Delta W_k$ can be replaced
by $\pm\sqrt h$ coin flips without changing the weak order \citep[Section~14.1]{kloeden-platen}.

\begin{theorem}[Milstein]\label{thm:ps-milstein}
Suppose (A1)--(A3) hold, that $a$ and $b$ are regular enough for the functions $b\,\partial_xb$,
$L^0a$, $L^1a$, $L^0b$, $L^0(b\,\partial_xb)$ and $L^1(b\,\partial_xb)$ to exist and be
continuous, where $L^0=\partial_t+a\,\partial_x+\frac12b^2\partial_{xx}$ and
$L^1=b\,\partial_x$, that $b\,\partial_xb$ is Lipschitz in $x$ uniformly in $t$, and that all
these functions grow at most linearly in $x$, uniformly in $t$.
\begin{enumerate}[label=(\roman*)]
  \item The Milstein scheme has strong order $1$: $\E|X_T-\tilde X^h_T|\le C\,h$
    \citep[Theorem~10.3.5, a case of Theorem~10.6.3]{kloeden-platen}.
  \item Under the hypotheses of \cref{thm:ps-euler}(ii), its weak order is $1$, and no
    better in general.
\end{enumerate}
The hypotheses hold, for example, when $a$ and $b$ are affine in $x$ with continuously
differentiable coefficients in $t$ (Brownian motion with drift, GBM, Ornstein--Uhlenbeck).
\end{theorem}

For (ii), the proof of \cref{thm:ps-euler}(ii) applies verbatim: the extra term
$\frac12b\,\partial_xb\,((\Delta W)^2-h)$ has mean zero and changes the second and third
moments of the step by $O(h^2)$, so the local weak error is still of order $h^2$, and it is not
smaller in general. Milstein improves the paths, not the order of the bias of prices.

For GBM, $\dd S=\mu S\,\dd t+\sigma S\,\dd W$, both schemes multiply by a factor at each step:
$1+\mu h+\sigma\Delta W_k$ for Euler, plus $\frac12\sigma^2((\Delta W_k)^2-h)$ for Milstein.
The exact factor is
$e^{(\mu-\sigma^2/2)h+\sigma\Delta W_k}=1+\mu h+\sigma\Delta W_k+\frac12\sigma^2((\Delta W_k)^2-h)+O(h^{3/2})$,
so the Milstein term is precisely the leading local error of Euler. Since GBM has an exact
solution on the same Brownian path, $S_0e^{(\mu-\sigma^2/2)T+\sigma W_T}$, it is the natural
laboratory for both orders.

\begin{figure}[ht]
  \centering
  \includegraphics[width=\textwidth]{ch07-orders}
  \caption{Euler and Milstein for GBM with the test-case parameters ($\mu=r-y=-1\%$,
    $\sigma=20\%$, $T=1$), $10^6$ paths, seed 2026. Left: strong error $\E|X_T-Y^h_T|$
    against the exact solution on the same path. Right: weak error
    $\E f(Y^h_T)-\E f(X_T)$ for $f(x)=5\log(1+e^{(x-100)/5})$, estimated from the
    differences on the same paths ($\pm2$ standard errors), and the Richardson-extrapolated
    Milstein error (absolute value) of \cref{prop:ps-richardson}; crosses: exact values.}
  \label{fig:ps-orders}
\end{figure}

\Cref{fig:ps-orders} shows the measurement with $m=1,2,4,\dots,64$ steps, all schemes driven
by the same $10^6$ Brownian paths (the coarse increments are sums of fine ones). At
$h=1/64$ the strong errors are $0.279$ for Euler and $0.0047$ for Milstein; the fitted slopes
over $m\ge4$ are $0.49$ and $0.99$. For the weak error we use the smooth call-like payoff
$f(x)=5\log(1+e^{(x-K)/5})$, $K=100$, and estimate $\E f(Y^h_T)-\E f(X_T)$ by the average of
$f(Y^h_T)-f(X_T)$ over the paths, which is unbiased for the difference and has a small
variance because the two terms are close. For GBM the weak errors can also be computed
without simulation, since $\log Y^h_T$ is a sum of $m$ \iid\ terms with an explicit law
(\texttt{gbm\_scheme\_weak\_error} convolves it numerically). The exact values are
$0.01059$, $0.00540$, $0.00273$ for Euler at $m=16$, $32$, $64$, and $0.00451$, $0.00221$,
$0.00110$ for Milstein: both of order $h$, as the theorems predict, and the Monte Carlo
estimates agree with them within $2.1$ standard errors (Euler's standard errors are $0.0005$,
$0.0004$, $0.0003$; Milstein's are ten times smaller, because its strong error is smaller). The bias is small next to the price $\E f(X_T)=8.206$:
at $\sigma=20\%$ a few steps per year are enough for this payoff.

\begin{remark}[When the theorems do not apply]\label{rem:ps-failures}
The hypotheses are not decoration.
\begin{itemize}
  \item \emph{Non-Lipschitz coefficients.} The square-root diffusion
    $\dd v=\kappa(\theta-v)\,\dd t+\xi\sqrt v\,\dd W$ of the Heston model violates (A1) at
    $v=0$, and the Euler step can make $v$ negative, after which $\sqrt v$ is undefined.
    Fixes such as full truncation and Andersen's scheme are in \cref{ch:19-stochastic-vol}.
  \item \emph{Superlinear growth.} For a large class of equations with superlinearly
    growing coefficients and noise acting from the start, for example
    $\dd X=-X^3\dd t+\dd W$, the moments of the Euler approximation at time $T$ diverge as
    $h\to0$, although those of $X_T$ are finite \citep{hutzenthaler-jentzen-kloeden}. (Without
    noise, Euler for $x'=-x^3$ started at $x_0$ stays bounded as soon as $hx_0^2\le2$; the
    noise is what occasionally throws the scheme far enough to explode.)
  \item \emph{Positivity.} The Euler factor $1+\mu h+\sigma\Delta W_k$ for GBM is negative
    with probability $N\big(-(1+\mu h)/(\sigma\sqrt h)\big)$: negligible for moderate $h$,
    fatal for a payoff such as $\log S_T$ (\exref{ex:ps-euler-moments}).
  \item \emph{Nonsmooth functionals.} A digital payoff is not even continuous. The expansion
    \eqref{eq:ps-tt}, hence weak order $1$, still holds for bounded measurable $f$ when the
    coefficients are smooth and the diffusion satisfies a nondegeneracy condition of
    H\"ormander type \citep{bally-talay}. For a barrier or a maximum the error is of order $h^{1/2}$, because
    the path is only inspected on the grid; see \cref{prop:bar-bgk} (for GBM), \cref{rem:bar-general} (Euler for general diffusions) and
    \exref{ex:ps-richardson-sqrt}.
\end{itemize}
\end{remark}

\demo{07-path-simulation.ipynb}{measure the strong and weak errors of Euler and Milstein for
GBM, for any volatility, drift and number of paths.}

\section{Richardson extrapolation}\label{sec:ps-richardson}

An error expansion such as \eqref{eq:ps-tt} can be used, not just bounded.

\begin{proposition}[Richardson extrapolation]\label{prop:ps-richardson}
Let $Y^h$ be a scheme and $f$ a payoff such that, for some $c\in\R$ and $0<p<q$,
\begin{equation}\label{eq:ps-expansion}
  \E f(Y^h_T)=\E f(X_T)+c\,h^p+O(h^q)\qquad(h\to0),
\end{equation}
and set $R^h=\big(2^pf(Y^h_T)-f(Y^{2h}_T)\big)/(2^p-1)$.
\begin{enumerate}[label=(\roman*)]
  \item $\E R^h=\E f(X_T)+O(h^q)$.
  \item If $Y^h_T$ and $Y^{2h}_T$ are computed on the same Brownian path, $f$ is Lipschitz
    and $\E|Y^h_T-X_T|^2\to0$, then $\Var R^h\to\Var f(X_T)$. If they are computed on
    independent paths, $\Var R^h\to\frac{4^p+1}{(2^p-1)^2}\Var f(X_T)$, which is
    $5\Var f(X_T)$ for $p=1$.
\end{enumerate}
\end{proposition}

\begin{proof}
(i) $\E R^h=\frac{2^p(\theta+ch^p)-(\theta+c2^ph^p)}{2^p-1}+O(h^q)=\theta+O(h^q)$, with
$\theta=\E f(X_T)$; the constant $c$ cancels without being known.

(ii) With a Lipschitz constant $L$ for $f$,
\[
  R^h-f(X_T)=\frac{2^p\big(f(Y^h_T)-f(X_T)\big)-\big(f(Y^{2h}_T)-f(X_T)\big)}{2^p-1},
\]
whose $L^2$ norm is at most
$\frac{L}{2^p-1}\big(2^p\|Y^h_T-X_T\|_2+\|Y^{2h}_T-X_T\|_2\big)\to0$. Convergence in $L^2$
implies convergence of variances. On independent paths,
$\Var R^h=\big(4^p\Var f(Y^h_T)+\Var f(Y^{2h}_T)\big)/(2^p-1)^2$, and each variance tends to
$\Var f(X_T)$ by the same $L^2$ convergence.
\end{proof}

The assumption is the expansion \eqref{eq:ps-expansion}, not merely an error bound: a bound
$|\text{error}|\le Ch^p$ allows an oscillating error that extrapolation can make worse. For
Euler, \eqref{eq:ps-expansion} with $p=1$, $q=2$ is \eqref{eq:ps-tt}; the mean-square
convergence in (ii) holds under the hypotheses of \cref{thm:ps-euler,thm:ps-milstein},
whose proofs bound $\E|X_T-Y^h_T|^2$. The coupling in (ii) is what makes extrapolation
cheap: the pair costs $m+m/2$ steps and has the variance of one path, while independent
pairs pay a factor $5$ in variance. The same coupling of a fine and a coarse path is the
building block of multilevel Monte Carlo (\cref{prop:ml-variance-decay}).

In the experiment of \cref{fig:ps-orders}, Richardson applied to Milstein has exact errors
$-0.000329$, $-0.000083$, $-0.000021$ at $m=16$, $32$, $64$, against $0.00451$, $0.00221$,
$0.00110$ before extrapolation: they fall by a factor $4$ per halving of $h$, so $q=2$, and
the Monte Carlo estimates of the figure follow them. For Euler the extrapolated errors are
$0.0162$, $0.0038$, $0.00091$, $0.00022$, $0.000055$ at $m=4,\dots,64$, also of order $h^2$;
for $m\ge16$ they are below the noise of $10^6$ paths. If the error expansion is
in powers of $h^{1/2}$, as for discretely monitored extremes, the exponent $p=1$ is wrong and
the extrapolation must use $p=\frac12$ (\exref{ex:ps-richardson-sqrt}).

\section{Time steps or paths?}\label{sec:ps-allocation}

A discretised Monte Carlo estimator has two errors: the bias of the scheme, which falls with
the step, and the statistical error, which falls with the number of paths. With a fixed
computing budget, more steps per path means fewer paths. The mean-square error trades the two
(\cref{prop:mc-mse}).

Let $\hat\theta=\frac1n\sum_{i=1}^nf(Y^h_{T,i})$ be computed from $n$ independent paths with
$m=T/h$ steps, each step costing $\kappa$, so that a budget $C$ buys $n=C/(\kappa m)$ paths
(we ignore integer rounding). With bias $b(h)=\E f(Y^h_T)-\E f(X_T)$ and variance
$v(h)=\Var f(Y^h_T)$,
\begin{equation}\label{eq:ps-mse}
  \mathrm{MSE}(h)=b(h)^2+\frac{v(h)}n=b(h)^2+\frac{v(h)\,\kappa T}{C\,h} .
\end{equation}

\begin{proposition}[Allocation of effort]\label{prop:ps-allocation}
Suppose that $b(h)=c\,h^\beta(1+o(1))$ with $c\ne0$, $\beta>0$, that $v(h)\to v>0$ as $h\to0$,
and that $|b|$ is bounded away from $0$ on $[\varepsilon,h_0]$ for every $\varepsilon>0$.
\begin{enumerate}[label=(\roman*)]
  \item For the model $c^2h^{2\beta}+v\kappa T/(Ch)$ of \eqref{eq:ps-mse}, the optimal step
    is
    \[
      h^*=\Big(\frac{v\,\kappa\,T}{2\beta\,c^2\,C}\Big)^{1/(2\beta+1)},
    \]
    so the number of steps grows like $C^{1/(2\beta+1)}$ and the number of paths like
    $C^{2\beta/(2\beta+1)}$. At the optimum the variance equals $2\beta$ times the squared
    bias, and $\mathrm{MSE}^*=(1+2\beta)\,c^2(h^*)^{2\beta}=M^*C^{-2\beta/(2\beta+1)}$ with
    $M^*=(1+2\beta)\,c^{2/(2\beta+1)}\big(v\kappa T/(2\beta)\big)^{2\beta/(2\beta+1)}$.
  \item For every choice of steps $h_C\in(0,h_0]$,
    $\liminf_{C\to\infty}C^{2\beta/(2\beta+1)}\,\mathrm{MSE}(h_C)\ge M^*$, and the choice
    $h_C=h^*$ attains this limit. The root-mean-square error therefore cannot fall faster
    than $C^{-\beta/(2\beta+1)}$.
\end{enumerate}
\end{proposition}

\begin{proof}
(i) Let $g(h)=c^2h^{2\beta}+A/h$ with $A=v\kappa T/C$. Its derivative
$g'(h)=h^{-2}\big(2\beta c^2h^{2\beta+1}-A\big)$ is negative for $h<h^*$ and positive for
$h>h^*$, so $h^*$ is the global minimiser. At $h^*$, $A/h=2\beta c^2h^{2\beta}$; this equation says that the variance term is $2\beta$ times the squared bias, so
$g(h^*)=(1+2\beta)c^2(h^*)^{2\beta}$, and substituting $h^*$ gives $M^*$. The paths number
$n=Ch^*/(\kappa T)\propto C^{1-1/(2\beta+1)}$.

(ii) Fix $\eta\in(0,1)$ and pass to a sequence $C_k\to\infty$ along which
$C^{2\beta/(2\beta+1)}\mathrm{MSE}(h_C)$ converges to its liminf. Either
$\liminf_kh_{C_k}>0$: then $h_{C_k}\ge\varepsilon$ for some $\varepsilon>0$ and all large $k$,
the mean-square error is at least $\inf_{[\varepsilon,h_0]}b^2>0$, and the liminf is
$+\infty$. Or a further subsequence has $h_{C_k}\to0$: along it, for $k$ large,
$b(h_C)^2\ge(1-\eta)c^2h_C^{2\beta}$ and $v(h_C)\ge(1-\eta)v$, so
$\mathrm{MSE}(h_C)\ge(1-\eta)\,g(h_C)\ge(1-\eta)\,g(h^*)=(1-\eta)M^*C^{-2\beta/(2\beta+1)}$.
As $\eta$ is arbitrary, the liminf is at least $M^*$. With $h_C=h^*\to0$, the same
expansions with $1+\eta$ give a matching upper bound.
\end{proof}

\Citet{duffie-glynn} derived this allocation rule; the mean-square version above follows
\textcite{glasserman-mc}, Section~1.1.3, which also allows a cost per path growing like
$h^{-\eta}$. The consequences, with the computing budget $C$:
\begin{center}
\begin{tabular}{lcc}
\toprule
method & bias order $\beta$ & root-mean-square error\\
\midrule
exact simulation (\cref{prop:ps-gbm-exact}) & none & $C^{-1/2}$\\
Euler or Milstein, optimal allocation & $1$ & $C^{-1/3}$\\
coupled Richardson extrapolation & $2$ & $C^{-2/5}$\\
multilevel Monte Carlo (\cref{thm:ml-complexity}) & --- & $C^{-1/2}$, up to a logarithm for Euler\\
\bottomrule
\end{tabular}
\end{center}
For a first-order scheme, the rule of thumb is $n\propto m^2$: doubling the number of steps
calls for four times as many paths. Spending the whole budget on paths with a fixed step
leaves the bias in place, and the error stalls at $|b(h)|$.

\begin{figure}[ht]
  \centering
  \includegraphics[width=\textwidth]{ch07-allocation}
  \caption{Allocation of effort for the payoff of \cref{fig:ps-orders} with Euler, using the
    constants $c=0.17$ and $v=154$; cost in units of one step of one path. Left: squared bias,
    variance and mean-square error against the number of steps for a budget of $10^6$
    (dots: exact squared biases). Right: root-mean-square error against the budget for exact
    simulation, Euler with the optimal allocation, and Euler with $4$ fixed steps (lines:
    model; dots: $100$ independent replications each).}
  \label{fig:ps-allocation}
\end{figure}

For the payoff of \cref{fig:ps-orders} the bias constant of Euler is $c\approx0.17$ ($m$
times the exact weak error is $0.169$ and $0.173$ at $m=16$ and $32$; the limit $h\to0$ is
$0.176$), and $v=\Var f(X_T)=154$. With a budget of
$C=10^6$ steps the optimal allocation is $m^*\approx7$ steps and $n^*\approx1.4\times10^5$
paths, for a root-mean-square error of $0.041$; exact simulation with $10^6$ paths at the same
cost reaches $0.012$. At $C=10^8$, $m^*\approx33$, $n^*\approx3\times10^6$, and the errors are
$0.0088$ against $0.0012$: the gap widens like $C^{1/6}$ (\cref{fig:ps-allocation}). When an
exact sampler exists, it should be used.

\demo{07-path-simulation.ipynb}{choose a budget and a bias constant and see the mean-square
error split between bias and variance as the number of steps changes.}

\furtherreading
\textcite{glasserman-mc} is the reference for this chapter: Section~3.1 for the three
constructions of Brownian paths, Section~3.2 for GBM, Sections~6.1--6.2 for the Euler and
Milstein schemes, convergence orders and extrapolation (Section~6.2.4), and Section~1.1.3 for
the allocation of effort. \textcite{kloeden-platen} prove the convergence theorems in full,
with the stochastic Taylor expansions behind every scheme (Chapter~5) and higher-order strong
and weak schemes (Chapters~10--15). \textcite{talay-tubaro} and \textcite{bally-talay} are the
sources for the weak error expansion, and \textcite{milstein-1974} introduced the scheme that
bears his name.
